% !PRINTED 399
\weberpage{399}
%<<CONT
(d. h. der Coëfficienten in ihrer Gleichung) ergiebt sich die, dass
das System, was man erhält, wenn man jede imaginäre Gerade
durch ihre conjugirte ersetzt, wieder ein Aronhold'sches System
derselben Curve ist. Denn dann ändern sich die Coëfficienten
nicht, wenn $i$ durch $-i$ ersetzt wird.

Die Curve wird also gewiss reell, wenn in dem gegebenen
Siebener-Systeme zu jeder imaginären Geraden die conjugirte
vorkommt; dann bilden sie ein \gesperrt{reelles Siebener-System}.

Ein solches reelles System kann nun enthalten:

1) eine reelle, drei Paar conjugirt imaginäre Geraden;

2) drei reelle und zwei Paar conjugirt imaginäre Geraden;

3) fünf reelle und ein Paar conjugirt imaginäre Geraden;

4) sieben reelle Geraden.

Wir werden sehen, dass diese vier Annahmen in derselben
Ordnung zu den im vorigen Paragraphen aufgezählten vier
Fällen führen.

Dieser Nachweis wird sehr einfach, wenn man sich der
Bezeichnungsweise der Doppeltangenten, die wir im \S.\,99 dargelegt
 haben, bedient, nach der wir die Steiner'schen Complexe
 unmittelbar aus der Bezeichnung bilden können.

1) Im ersten Falle bezeichnen wir die gegebenen sieben
Geraden mit
\[
0,\ 1,\ 1',\ 2,\ 2',\ 3,\ 3',
\]
setzen $0$ als reell voraus, $1$ mit $1'$, $2$ mit $2'$, $3$ mit $3'$ conjugirt
imaginär. Die übrigen 21 Doppeltangenten werden dann durch
die Zeichen $[01], [01'], [11'], \dots$ bezeichnet.

Nach der Vorschrift des \S.\,99 erhalten wir die beiden
folgenden Steiner'schen Complexe:
\[
\begin{aligned}
(1)&\quad 0[01],\ 1'[11'],\ 2[12],\ 2'[12'],\ 3[13],\ 3'[13']\\
(1')&\quad 0[01'],\ 1[11'],\ 2[1'2],\ 2'[1'2'],\ 3[1'3],\ 3'[1'3'].
\end{aligned}
\]

Der zu (1) conjugirte Complex muss die Elemente $0,1,2,2',3,3'$
enthalten, und ist also nach \S.\,98, 7. mit (1') identisch, der
diese sechs Elemente gleichfalls enthält.

Daraus folgt, dass $[11']$ reell ist, und dass
\[
\begin{aligned}
&{}[01]\quad [01']\\
&{}[12]\quad [1'2']\\
&{}[12']\quad [1'2]
\end{aligned}
\]
%CONT>>
